\appendix
\section{Apéndice A. Material gráfico complementario}

Las figuras de este apéndice documentan el análisis exploratorio de las series de tráfico. Se
excluyeron del cuerpo por resultar redundantes respecto de figuras allí incluidas, o por
referir a aspectos que el trabajo no modela.

\begin{figure}[htbp]
\centering
\includegraphics[width=0.85\textwidth]{alb_serie.png}
\caption{Serie horaria del balanceador de carga, análisis exploratorio}
\end{figure}

\begin{figure}[htbp]
\centering
\includegraphics[width=0.85\textwidth]{pos_serie.png}
\caption{Serie horaria de punto de venta, análisis exploratorio}
\end{figure}

\begin{figure}[htbp]
\centering
\includegraphics[width=0.85\textwidth]{alb_diaria.png}
\caption{Agregación diaria de la serie del balanceador de carga}
\end{figure}

\begin{figure}[htbp]
\centering
\includegraphics[width=0.85\textwidth]{pos_diaria.png}
\caption{Agregación diaria de la serie de punto de venta}
\end{figure}

\begin{figure}[htbp]
\centering
\includegraphics[width=0.85\textwidth]{alb_diasem.png}
\caption{Patrón por día de la semana, serie del balanceador de carga}
\end{figure}

\begin{figure}[htbp]
\centering
\includegraphics[width=0.85\textwidth]{pos_diasem.png}
\caption{Patrón por día de la semana, serie de punto de venta}
\end{figure}

\begin{figure}[htbp]
\centering
\includegraphics[width=0.85\textwidth]{alb_moviles.png}
\caption{Media y desvío estándar móviles, serie del balanceador de carga}
\end{figure}

\begin{figure}[htbp]
\centering
\includegraphics[width=0.85\textwidth]{pos_moviles.png}
\caption{Media y desvío estándar móviles, serie de punto de venta}
\end{figure}

\begin{figure}[htbp]
\centering
\includegraphics[width=0.85\textwidth]{alb_descomp.png}
\caption{Descomposición estacional de la serie del balanceador de carga}
\end{figure}

\begin{figure}[htbp]
\centering
\includegraphics[width=0.85\textwidth]{pos_descomp.png}
\caption{Descomposición estacional de la serie de punto de venta}
\end{figure}

\begin{figure}[htbp]
\centering
\includegraphics[width=0.85\textwidth]{trafico_traintest_1.png}
\caption{Ajuste sobre el conjunto de prueba, serie del balanceador de carga}
\end{figure}

\begin{figure}[htbp]
\centering
\includegraphics[width=0.85\textwidth]{trafico_traintest_2.png}
\caption{Ajuste sobre el conjunto de prueba, serie de punto de venta}
\end{figure}

\section{Apéndice B. Rutinas de cómputo}

\subsection*{Construcción de la serie de altitud}

La rutina siguiente extrae la altitud del registro binario de telemetría, la lleva a una grilla
regular de un segundo y aísla el tramo contiguo más extenso en vuelo.

\begin{lstlisting}
def leer_altitud(ruta):
    """Extrae la altitud relativa, en metros, de un registro MAVLink."""
    conexion = mavutil.mavlink_connection(ruta)
    registros = []
    while True:
        mensaje = conexion.recv_match(type="GLOBAL_POSITION_INT")
        if mensaje is None:
            break
        registros.append((mensaje._timestamp, mensaje.relative_alt / 1000.0))
    crudo = pd.DataFrame(registros, columns=["timestamp", "altitud_m"])
    crudo["timestamp"] = pd.to_datetime(crudo["timestamp"], unit="s")
    return crudo.set_index("timestamp")["altitud_m"]


def bloque_contiguo_mas_largo(mascara):
    """Indices de inicio y fin del tramo contiguo mas largo que cumple la mascara."""
    mejor_inicio, mejor_largo = 0, 0
    inicio = None
    for posicion, activo in enumerate(mascara):
        if activo and inicio is None:
            inicio = posicion
        elif not activo and inicio is not None:
            if posicion - inicio > mejor_largo:
                mejor_inicio, mejor_largo = inicio, posicion - inicio
            inicio = None
    if inicio is not None and len(mascara) - inicio > mejor_largo:
        mejor_inicio, mejor_largo = inicio, len(mascara) - inicio
    return mejor_inicio, mejor_inicio + mejor_largo
\end{lstlisting}

\subsection*{Representación conjunta de las tres funciones}

\begin{lstlisting}
def graficar_fas_fac_facp(datos, titulo, rezagos=40):
    """Grafica autocovarianza, autocorrelacion y parcial en una sola figura."""
    figura, ejes = plt.subplots(1, 3, figsize=(13, 3.4))
    autocovarianza = sm.tsa.acovf(datos, nlag=rezagos, fft=True)
    ejes[0].stem(range(len(autocovarianza)), autocovarianza, basefmt=" ")
    ejes[0].set_title("FAS - autocovarianza")
    plot_acf(datos, lags=rezagos, ax=ejes[1], title="FAC - autocorrelacion")
    plot_pacf(datos, lags=rezagos, ax=ejes[2], method="ywm",
              title="FACP - autocorrelacion parcial")
    figura.suptitle(titulo, y=1.05)
    plt.tight_layout()
    plt.show()
\end{lstlisting}

\subsection*{Comparación de especificaciones del modelo multivariado}

La rutina siguiente evalúa cada orden de rezagos e invierte la transformación antes de calcular
las métricas, de modo que el error quede expresado en la escala original.

\begin{lstlisting}
def evaluar(p, entrenamiento, prueba, log_original, alineadas, horizonte):
    resultado = VAR(entrenamiento).fit(p)
    pronostico = pd.DataFrame(
        resultado.forecast(entrenamiento.values[-p:], horizonte),
        index=prueba.index, columns=prueba.columns)
    filas = []
    for columna in prueba.columns:
        historico = log_original[columna].copy()
        prediccion = []
        for i, t in enumerate(prueba.index):
            base = historico.loc[t - pd.Timedelta(hours=24)]
            valor = base + pronostico[columna].iloc[i]
            historico.loc[t] = valor
            prediccion.append(valor)
        estimado = np.expm1(np.array(prediccion))
        real = alineadas[columna].reindex(prueba.index).to_numpy()
        filas.append({
            "rezagos": p,
            "serie": columna,
            "MAE": np.mean(np.abs(estimado - real)),
            "RMSE": np.sqrt(np.mean((estimado - real) ** 2)),
            "MAPE": np.mean(np.abs((estimado - real) / real)) * 100,
        })
    return filas
\end{lstlisting}
